% unit: u01
% sources: sources/069-the-geometry-of-linear-equations.txt, sources/115-the-geometry-of-linear-equations.txt, sources/116-mit18-06scf11-ses1-1sum-pdf-download-file.txt, sources/117-geometry-of-linear-algebra.txt, sources/118-mit18-06scf11-ses1-1prob-pdf-download-file.txt, sources/119-mit18-06scf11-ses1-1sol-pdf-download-file.txt, short_plan.json (s01)
\section{The geometry of linear equations}\label{sec:u01}

The central problem of linear algebra is to solve a system of linear equations, and the
course starts with the friendliest case: $n$ equations in $n$ unknowns. This section looks
at one small system in three ways. The \emph{row picture} draws one equation at a time and
looks for the point where the graphs meet. The \emph{column picture} reads the same system
as a question about combining vectors. The \emph{matrix form} packs everything into the
single equation $A\vb{x}=\vb{b}$. The column picture is the one that survives in ten or a
thousand dimensions, and it leads directly to the question that drives the first part of the
course: for which right-hand sides $\vb{b}$ can the system be solved at all?

\subsection{One system, three readings}

Our first example is the system
\begin{equation}\label{eq:u01-system}
\begin{aligned}
  2x - y &= 0,\\
  -x + 2y &= 3 .
\end{aligned}
\end{equation}
All the information sits in the numbers multiplying the unknowns and in the numbers on the
right, so we collect them.

\begin{definition}[Matrix, coefficient matrix, $A\vb{x}=\vb{b}$]\label{def:u01-matrix}
A \emph{matrix} is a rectangular array of numbers. An $m\times n$ matrix has $m$ rows and
$n$ columns, and $a_{ij}$ denotes the entry in row $i$ and column $j$. For a system of $m$
linear equations in $n$ unknowns $x_1,\dots,x_n$, the \emph{coefficient matrix} $A$ holds
the coefficient of $x_j$ in equation $i$ as its entry $a_{ij}$; the unknowns form the vector
$\vb{x}=(x_1,\dots,x_n)$ and the right-hand sides form the vector $\vb{b}=(b_1,\dots,b_m)$.
The system is written $A\vb{x}=\vb{b}$.
\end{definition}

Vectors are written as columns; to save space we also write them as lists in parentheses,
so $\vb{b}=(0,3)$ means the column with entries $0$ and $3$. For the system
\eqref{eq:u01-system},
\[
  A=\begin{bmatrix} 2 & -1\\ -1 & 2\end{bmatrix},\qquad
  \vb{x}=\begin{bmatrix} x\\ y\end{bmatrix},\qquad
  \vb{b}=\begin{bmatrix} 0\\ 3\end{bmatrix},\qquad
  \begin{bmatrix} 2 & -1\\ -1 & 2\end{bmatrix}\begin{bmatrix} x\\ y\end{bmatrix}
  =\begin{bmatrix} 0\\ 3\end{bmatrix}.
\]
At this point $A\vb{x}$ is shorthand for ``the left-hand sides''. In
Section~\ref{subsec:u01-product} we make it a genuine product of a matrix and a vector.

\subsection{The row picture: lines and planes that meet}

In the row picture each equation is drawn separately in the $xy$-plane. All points
satisfying $2x-y=0$ lie on a straight line (the word \emph{linear} comes from \emph{line}).
It passes through the origin, because $(0,0)$ satisfies it, and through $(1,2)$ and
$(\tfrac12,1)$. The second line $-x+2y=3$ does \emph{not} pass through the origin, since
$x=y=0$ gives $0\neq 3$; two of its points are $(-3,0)$ and $(-1,1)$. A solution of the
system must lie on both lines, so it is their crossing point. The left panel of
Figure~\ref{fig:u01-pictures} shows the lines meeting at $(x,y)=(1,2)$, and substituting
confirms it: $2\cdot1-2=0$ and $-1+2\cdot2=3$.

Two lines in a plane can do three things. They can cross at one point (one solution), they
can be parallel and distinct (no solution), or they can coincide (infinitely many solutions).
The nice case $n$ equations, $n$ unknowns, one solution, is the first; the other two are the
singular cases that the course will learn to detect.

In three unknowns each equation describes a \emph{plane} in $xyz$-space. Take
\begin{equation}\label{eq:u01-three}
\begin{aligned}
  2x - y \phantom{{}-z} &= 0,\\
  -x + 2y - z &= -1,\\
  -3y + 4z &= 4 .
\end{aligned}
\end{equation}
The second equation is not satisfied by the origin; it contains $(1,0,0)$, $(0,0,1)$ and
$(0,-\tfrac12,0)$, and those three points determine its plane. The first equation is another
plane (it does not involve $z$, so it is ``vertical''). Two non-parallel planes meet in a
line, and the third plane cuts that line in a single point, the solution. The picture is
already hard to draw for three planes, and in four dimensions it is hopeless. We need a
better view.

\subsection{The column picture: combinations of columns}

Read the system \eqref{eq:u01-system} a column at a time. The unknown $x$ multiplies the
first column of $A$ and $y$ multiplies the second:
\begin{equation}\label{eq:u01-columns}
  x\begin{bmatrix} 2\\ -1\end{bmatrix} + y\begin{bmatrix} -1\\ 2\end{bmatrix}
  = \begin{bmatrix} 0\\ 3\end{bmatrix}.
\end{equation}
Vectors are added component by component and multiplied by a number (a \emph{scalar}) in
every component, so \eqref{eq:u01-columns} is exactly the two equations again. Now the question
is different: which amounts of the two column vectors add up to the vector $\vb{b}$?

\begin{definition}[Linear combination]\label{def:u01-combination}
Given vectors $\vb{v}_1,\dots,\vb{v}_n$ in $\R^m$ and scalars $c_1,\dots,c_n$, the vector
\[
  c_1\vb{v}_1 + c_2\vb{v}_2 + \cdots + c_n\vb{v}_n
\]
is a \emph{linear combination} of $\vb{v}_1,\dots,\vb{v}_n$. The scalars $c_j$ are its
\emph{weights} (or coefficients).
\end{definition}

Taking linear combinations is the basic operation of the whole subject. The column picture
says: solving $A\vb{x}=\vb{b}$ means finding the linear combination of the columns of $A$ that
produces $\vb{b}$, and the weights in that combination are the unknowns.

\begin{example}[Two pictures of one system]\label{ex:u01-two-pictures}
Solve \eqref{eq:u01-system} in both pictures. The columns of $A$ are
$\vb{a}_1=(2,-1)$ and $\vb{a}_2=(-1,2)$, and the target is $\vb{b}=(0,3)$.
\begin{enumerate}
  \item \emph{Row picture.} The lines $y=2x$ and $-x+2y=3$ cross where
    $-x+2(2x)=3$, that is $3x=3$. So $x=1$, $y=2$, and indeed $-1+2\cdot 2=3$.
  \item \emph{Column picture.} Take $x=1$ copy of $\vb{a}_1$: the arrow from the origin to
    $(2,-1)$.
  \item Take $y=2$ copies of $\vb{a}_2$: $2(-1,2)=(-2,4)$.
  \item Add tip to tail: $(2,-1)+(-2,4)=(2-2,\,-1+4)=(0,3)$, which is exactly $\vb{b}$.
\end{enumerate}
The same two numbers $x=1$, $y=2$ are the crossing point of the rows and the recipe
$1\,\vb{a}_1+2\,\vb{a}_2=\vb{b}$ for the columns. Figure~\ref{fig:u01-pictures} shows both.
\end{example}

\begin{figure}[ht]
\centering
\begin{tikzpicture}[scale=0.72,>=Stealth]
  \draw[->] (-3.4,0) -- (3.2,0) node[right] {$x$};
  \draw[->] (0,-1.4) -- (0,4.5) node[above] {$y$};
  \foreach \t in {-3,-2,-1,1,2} \draw (\t,0.08) -- (\t,-0.08) node[below] {\scriptsize $\t$};
  \foreach \t in {1,2,3,4} \draw (0.08,\t) -- (-0.08,\t) node[left] {\scriptsize $\t$};
  \draw[very thick,blue] (-0.65,-1.3) -- (2.2,4.4) node[right] {$2x-y=0$};
  \draw[very thick,orange!85!black] (-3.4,-0.2) -- (3.2,3.1) node[below right] {$-x+2y=3$};
  \fill (1,2) circle (2.5pt) node[below right=1pt] {$(1,2)$};
  \node[font=\small] at (0,-2.2) {row picture: $(x,y)$-plane};
\end{tikzpicture}
\qquad
\begin{tikzpicture}[scale=0.72,>=Stealth]
  \draw[->] (-2.8,0) -- (3,0) node[right] {$b_1$};
  \draw[->] (0,-1.4) -- (0,4.5) node[above] {$b_2$};
  \foreach \t in {-2,-1,1,2} \draw (\t,0.08) -- (\t,-0.08) node[below] {\scriptsize $\t$};
  \foreach \t in {1,2,3,4} \draw (0.08,\t) -- (-0.08,\t) node[left] {\scriptsize $\t$};
  \draw[->,very thick,blue] (0,0) -- (2,-1) node[below right] {$\vb{a}_1=(2,-1)$};
  \draw[->,very thick,orange!85!black] (0,0) -- (-1,2) node[left] {$\vb{a}_2$};
  \draw[->,thick,dashed,orange!85!black] (2,-1) -- (1,1);
  \draw[->,thick,dashed,orange!85!black] (1,1) -- (0,3);
  \node[orange!85!black,right,font=\small] at (1.05,1.5) {$2\vb{a}_2=(-2,4)$};
  \draw[->,very thick] (0,0) -- (0,3);
  \fill (0,3) circle (2.5pt) node[above right] {$\vb{b}=(0,3)$};
  \node[font=\small] at (0,-2.2) {column picture: $1\,\vb{a}_1+2\,\vb{a}_2=\vb{b}$};
\end{tikzpicture}
\caption{The system $2x-y=0$, $-x+2y=3$ two ways. Left: each equation is a line and the
solution is the crossing point $(1,2)$. Right: the columns $\vb{a}_1,\vb{a}_2$ are arrows in
the plane of right-hand sides; one $\vb{a}_1$ followed by two copies of $\vb{a}_2$ (dashed,
tip to tail) lands on $\vb{b}=(0,3)$. The weights $1,2$ are the coordinates of the crossing
point.}\label{fig:u01-pictures}
\end{figure}

The two panels live in different planes. In the row picture the axes are the unknowns
$x$ and $y$, and the answer is a point. In the column picture the axes are the components
$b_1,b_2$ of the right-hand side (the first and second equations), and the answer is a pair
of weights.

What do \emph{all} combinations $x\vb{a}_1+y\vb{a}_2$ give? Multiples of $\vb{a}_1$ fill a
line through the origin, multiples of $\vb{a}_2$ fill another line, and because the two
lines point in different directions, adding a point of one to a point of the other reaches
every point of the plane. So for this $A$ the system can be solved for \emph{every}
right-hand side $\vb{b}$ (Exercise~\ref{ex:u01-every-b} finds the weights explicitly).

The column picture of the $3\times3$ system \eqref{eq:u01-three} is
\[
  x\begin{bmatrix} 2\\ -1\\ 0\end{bmatrix}
  + y\begin{bmatrix} -1\\ 2\\ -3\end{bmatrix}
  + z\begin{bmatrix} 0\\ -1\\ 4\end{bmatrix}
  = \begin{bmatrix} 0\\ -1\\ 4\end{bmatrix}.
\]
This right-hand side was chosen to equal the third column, so the combination is visible at
once: $x=0$, $y=0$, $z=1$. That is the point where the three planes of the row picture meet,
found without drawing a single plane. If we change the right-hand side to
$\vb{b}=(1,1,-3)$, which is column~1 plus column~2, the solution becomes $(x,y,z)=(1,1,0)$.
In the row picture all three planes moved; in the column picture the three columns stayed
put and only the target changed. Usually the combination is not visible by eye, and the
systematic method of finding it, elimination, is the subject of Section~\ref{sec:u03}.

\subsection{Multiplying a matrix by a vector}\label{subsec:u01-product}

The column picture tells us what $A\vb{x}$ should mean.

\begin{definition}[Matrix times vector]\label{def:u01-Ax}
Let $A$ be an $m\times n$ matrix with columns $\vb{a}_1,\dots,\vb{a}_n$ (each in $\R^m$),
and let $\vb{x}=(x_1,\dots,x_n)\in\R^n$. The product $A\vb{x}$ is the combination of the
columns with weights from $\vb{x}$:
\[
  A\vb{x} = x_1\vb{a}_1 + x_2\vb{a}_2 + \cdots + x_n\vb{a}_n \in \R^m .
\]
\end{definition}

The product is defined only when the number of columns of $A$ equals the number of
components of $\vb{x}$. The familiar ``row times vector'' rule is the same computation
organised differently.

\begin{proposition}[Row form of $A\vb{x}$]\label{prop:u01-rows}
For $A$ and $\vb{x}$ as in Definition~\ref{def:u01-Ax}, component $i$ of $A\vb{x}$ is the
dot product of row $i$ of $A$ with $\vb{x}$:
\[
  (A\vb{x})_i = a_{i1}x_1 + a_{i2}x_2 + \cdots + a_{in}x_n = \sum_{j=1}^{n} a_{ij}x_j .
\]
\end{proposition}

\begin{proof}
Component $i$ of the vector $x_j\vb{a}_j$ is $x_j$ times component $i$ of $\vb{a}_j$, which
is $a_{ij}x_j$. Components of a sum are sums of components, so component $i$ of
$x_1\vb{a}_1+\cdots+x_n\vb{a}_n$ is $a_{i1}x_1+\cdots+a_{in}x_n$, and this is row $i$ of $A$
dotted with $\vb{x}$.
\end{proof}

So the row picture corresponds to computing $A\vb{x}$ by rows (each equation is one dot
product), and the column picture corresponds to computing $A\vb{x}$ by columns. With
Definition~\ref{def:u01-Ax}, the shorthand $A\vb{x}=\vb{b}$ for a linear system is now a true
equation between vectors.

\begin{example}[Both ways]\label{ex:u01-both-ways}
Compute $\begin{bmatrix} 2 & 5\\ 1 & 3\end{bmatrix}\begin{bmatrix} 1\\ 2\end{bmatrix}$.
By columns, take one of the first column and two of the second:
\[
  1\begin{bmatrix} 2\\ 1\end{bmatrix} + 2\begin{bmatrix} 5\\ 3\end{bmatrix}
  = \begin{bmatrix} 2+10\\ 1+6\end{bmatrix} = \begin{bmatrix} 12\\ 7\end{bmatrix}.
\]
By rows, $(2,5)\cdot(1,2) = 2+10 = 12$ and $(1,3)\cdot(1,2)=1+6=7$. The same four
multiplications happen in a different order. In particular, if the right-hand side were
$\vb{b}=(12,7)$, the solution of $A\vb{x}=\vb{b}$ would be $\vb{x}=(1,2)$.
\end{example}

\subsection{Can every right-hand side be reached?}

For a given square matrix $A$, can we solve $A\vb{x}=\vb{b}$ for \emph{every} $\vb{b}$?
In the language of the column picture the question becomes: do the linear combinations of
the columns fill the whole space $\R^n$? For the $2\times2$ example the answer was yes,
and for the $3\times3$ matrix of \eqref{eq:u01-three} it is also yes. Such matrices are called
\emph{nonsingular} or \emph{invertible}; Section~\ref{sec:u04} defines the inverse matrix
that undoes $A$. When the answer is no, $A$ is \emph{singular}.

Here is how the answer can fail in three dimensions. If the three columns happen to lie in
one plane through the origin, every combination of them lies in that plane too, and no
right-hand side off the plane can be produced.

\begin{proposition}[A dependent column adds nothing]\label{prop:u01-dependent}
Let $A$ have columns $\vb{a}_1,\vb{a}_2,\vb{a}_3$ in $\R^3$ and suppose
$\vb{a}_3 = c_1\vb{a}_1 + c_2\vb{a}_2$. Then every vector $A\vb{x}$ is a combination of
$\vb{a}_1$ and $\vb{a}_2$ alone.
\end{proposition}

\begin{proof}
For any $\vb{x}=(x_1,x_2,x_3)$,
\[
  A\vb{x} = x_1\vb{a}_1 + x_2\vb{a}_2 + x_3(c_1\vb{a}_1+c_2\vb{a}_2)
  = (x_1+c_1x_3)\,\vb{a}_1 + (x_2+c_2x_3)\,\vb{a}_2 . \qedhere
\]
\end{proof}

The combinations of two vectors in $\R^3$ fill at most a plane through the origin, so a
matrix with a dependent column is singular: some $\vb{b}$ are out of reach. (That two
vectors cannot fill $\R^3$ is intuitively clear; the precise statement, that a space of
dimension $3$ needs three independent vectors, is proved in Section~\ref{sec:u10}.)

\begin{example}[Three columns in a plane]\label{ex:u01-plane}
Let $\vb{w}_1=(1,2,3)$, $\vb{w}_2=(4,5,6)$, $\vb{w}_3=(7,8,9)$. Find a combination
$x_1\vb{w}_1+x_2\vb{w}_2+x_3\vb{w}_3$ that gives the zero vector, and a vector $\vb{b}$ that
is not a combination of them.

The requirement is three equations:
\[
  x_1+4x_2+7x_3=0,\qquad 2x_1+5x_2+8x_3=0,\qquad 3x_1+6x_2+9x_3=0 .
\]
Subtract twice the first equation from the second: $-3x_2-6x_3=0$. Subtract three times
the first from the third: $-6x_2-12x_3=0$, which is twice the previous equation and so
tells us nothing new. Choose $x_3=1$; then $x_2=-2$ and $x_1=-4x_2-7x_3=8-7=1$. Hence
\[
  \vb{w}_1 - 2\vb{w}_2 + \vb{w}_3 = (1-8+7,\; 2-10+8,\; 3-12+9) = (0,0,0).
\]
The vectors are \emph{dependent}: $\vb{w}_2=\tfrac12(\vb{w}_1+\vb{w}_3)$ lies in the plane
of $\vb{w}_1$ and $\vb{w}_3$, so all three lie in one plane through the origin. Which plane?
Each $\vb{w}_j=(w_1,w_2,w_3)$ satisfies $w_1-2w_2+w_3=0$ ($1-4+3=0$, $4-10+6=0$,
$7-16+9=0$), and that equation is preserved by adding vectors and multiplying by scalars. So
every combination $\vb{b}$ satisfies $b_1-2b_2+b_3=0$. The vector $\vb{b}=(1,0,0)$ gives
$1\neq0$, so $W\vb{x}=(1,0,0)$ has no solution for the matrix $W$ with columns
$\vb{w}_1,\vb{w}_2,\vb{w}_3$, and $W$ is not invertible.
\end{example}

The same thinking works in any dimension, even where nothing can be drawn. Nine equations
in nine unknowns give nine columns, each a vector in $\R^9$. A matrix with random entries
almost certainly has columns whose combinations fill all of $\R^9$. But if, say, the ninth
column equals the eighth, it contributes nothing new, and the combinations fill only an
eight-dimensional ``plane'' through the origin inside $\R^9$. Learning to think about such
flat pieces of high-dimensional space is the heart of the course.

\begin{trap}
Do not mix up the two pictures. In the row picture the point $(1,2)$ is a \emph{location}
in the plane of unknowns; in the column picture the same numbers are \emph{weights}, and the
arrows live in the plane of right-hand sides, where the target is $\vb{b}=(0,3)$, not
$(1,2)$. Likewise, $A\vb{x}$ is not only ``rows dotted with $\vb{x}$'': it is a combination
of the columns of $A$, and that is the view that keeps working when no lines or planes can
be drawn.
\end{trap}

\subsection{Exercises}

\begin{exercise}\label{ex:u01-multiply}
Compute $\begin{bmatrix} 1 & 2 & 0\\ 2 & 0 & 3\\ 4 & 1 & 1\end{bmatrix}
\begin{bmatrix} 3\\ -2\\ 1\end{bmatrix}$ by rows, and again as a combination of columns.
\end{exercise}
\begin{solution}
By rows: $1\cdot3+2\cdot(-2)+0\cdot1=-1$, then $2\cdot3+0\cdot(-2)+3\cdot1=9$, then
$4\cdot3+1\cdot(-2)+1\cdot1=11$. By columns:
\[
  3\begin{bmatrix} 1\\ 2\\ 4\end{bmatrix} - 2\begin{bmatrix} 2\\ 0\\ 1\end{bmatrix}
  + 1\begin{bmatrix} 0\\ 3\\ 1\end{bmatrix}
  = \begin{bmatrix} 3-4+0\\ 6-0+3\\ 12-2+1\end{bmatrix}
  = \begin{bmatrix} -1\\ 9\\ 11\end{bmatrix}. \qedhere
\]
\end{solution}

\begin{exercise}\label{ex:u01-shapes}
True or false: a $3\times2$ matrix $A$ times a $2\times3$ matrix $B$ is a $3\times3$ matrix
$AB$. State the general rule for the shapes. (Think of $AB$ as $A$ multiplying each column
of $B$; Section~\ref{sec:u04} develops this.)
\end{exercise}
\begin{solution}
True. Each column of $B$ has $2$ components, which matches the $2$ columns of $A$, so $A$
can multiply it and produces a vector with $3$ components. $B$ has $3$ columns, so $AB$ has
$3$ columns of length $3$. In general an $m\times n$ matrix times an $n\times p$ matrix is
$m\times p$: the inner dimensions must agree and the outer ones give the shape.
\end{solution}

\begin{exercise}\label{ex:u01-recitation}
Solve $2x+y=3$, $x-2y=-1$. Sketch the row picture (find two points on each line) and the
column picture, and write the system in matrix form.
\end{exercise}
\begin{solution}
From the second equation $x=2y-1$; substituting, $2(2y-1)+y=3$, so $5y=5$, $y=1$ and $x=1$.
Row picture: $2x+y=3$ passes through $(0,3)$ and $(1,1)$; $x-2y=-1$ passes through
$(0,\tfrac12)$ and $(1,1)$; the lines cross at $(1,1)$. Column picture: the columns are
$\vb{v}_1=(2,1)$ and $\vb{v}_2=(1,-2)$, and one of each gives the diagonal of the
parallelogram they span, $(2+1,\,1-2)=(3,-1)=\vb{b}$. Matrix form:
$\begin{bmatrix} 2 & 1\\ 1 & -2\end{bmatrix}\begin{bmatrix} x\\ y\end{bmatrix}
=\begin{bmatrix} 3\\ -1\end{bmatrix}$.
\end{solution}

\begin{exercise}\label{ex:u01-every-b}
With the columns $\vb{a}_1=(2,-1)$ and $\vb{a}_2=(-1,2)$ of Example~\ref{ex:u01-two-pictures},
which combination reaches $\vb{b}=(1,1)$? Then find the weights $x,y$ with
$x\vb{a}_1+y\vb{a}_2=(b_1,b_2)$ for an arbitrary right-hand side, confirming that every $\vb{b}$
is reachable.
\end{exercise}
\begin{solution}
$\vb{a}_1+\vb{a}_2=(2-1,\,-1+2)=(1,1)$, so $x=y=1$. In general we need $2x-y=b_1$ and
$-x+2y=b_2$. Adding twice the first equation to the second gives $3x=2b_1+b_2$; adding the
first to twice the second gives $3y=b_1+2b_2$. So
\[
  x=\frac{2b_1+b_2}{3},\qquad y=\frac{b_1+2b_2}{3}.
\]
Check: $2x-y=\frac{4b_1+2b_2-b_1-2b_2}{3}=b_1$ and $-x+2y=\frac{-2b_1-b_2+2b_1+4b_2}{3}=b_2$.
Weights exist for every $\vb{b}$, so the columns fill the plane. For $\vb{b}=(1,1)$ the
formula gives $x=y=1$ again.
\end{solution}

\begin{exercise}\label{ex:u01-param}
The matrix $A$ has columns $(1,0,1)$, $(0,1,1)$ and $(1,1,c)$. For $c=2$, show that the
columns lie in a plane, find an equation of that plane, and give a right-hand side $\vb{b}$
for which $A\vb{x}=\vb{b}$ has no solution. For $c=3$, solve $A\vb{x}=(0,0,1)$.
\end{exercise}
\begin{solution}
For $c=2$ the third column is $(1,0,1)+(0,1,1)$, so by
Proposition~\ref{prop:u01-dependent} every $A\vb{x}$ is a combination
$s(1,0,1)+t(0,1,1)=(s,\,t,\,s+t)$. These vectors satisfy $b_3=b_1+b_2$, the equation of
the plane. The vector $\vb{b}=(0,0,1)$ has $b_3=1\neq0=b_1+b_2$, so it is unreachable.

For $c=3$ we need $x_1+x_3=0$, $x_2+x_3=0$ and $x_1+x_2+3x_3=1$. The first two give
$x_1=x_2=-x_3$, and the third becomes $-x_3-x_3+3x_3=x_3=1$. So $\vb{x}=(-1,-1,1)$, and
indeed $-(1,0,1)-(0,1,1)+(1,1,3)=(0,0,1)$. The column $(1,1,3)$ is off the plane
$b_3=b_1+b_2$, which is why the target off that plane can now be reached.
\end{solution}
